% ============================================================
\section{Individual Contributions}
% ============================================================

\textbf{Rahul Ranjan Sah (50\%)} extended Module~1 with research-backed distance metrics---Krumhansl-Kessler key profiles, circle-of-fifths distance, and Bhattacharyya distance for MFCC similarity---on top of the original architecture. For Module~2, he developed the bidirectional beam search algorithm and implemented all MetaBrainz API clients: the MusicBrainzDB PostgreSQL client with connection pooling, the ListenBrainz client, and the AcousticBrainz client. He also built the six constraint types and min-conflicts solver in Module~3, implemented user preference learning via exponential moving average, integrated the Essentia fallback for tracks missing AcousticBrainz data, conducted performance benchmarking, wrote unit tests for Modules 1 and 2, and managed research review and citation gathering.

\textbf{Michael Thomas (50\%)} designed the original Module~1 architecture and ProbLog integration. He was responsible for all data sourcing infrastructure: deploying the self-hosted MusicBrainz PostgreSQL mirror and building the canonical song search API on top of it. For Module~3, he built the playlist assembly pipeline and the three-level explanation system. For Module~4, he created the feature engineering pipeline, trained the logistic regression and MLP classifiers, and built the training data pipeline with database integration. He also implemented the Flask REST API and developed the React frontend, and wrote unit tests for Modules 3 and 4.

Both team members participated in architecture design sessions, conducted code reviews via pull requests, performed integration testing collaboratively, and shared documentation responsibilities across modules.
